% ============================================================
%  Preamble – Packages & Layout
%  Eingebunden in main.tex via % ============================================================
%  Preamble – Packages & Layout
%  Eingebunden in main.tex via % ============================================================
%  Preamble – Packages & Layout
%  Eingebunden in main.tex via % ============================================================
%  Preamble – Packages & Layout
%  Eingebunden in main.tex via \input{preamble}
% ============================================================

% ── Encoding & Language ──────────────────────────────────────
\usepackage[utf8]{inputenc}
\usepackage[T1]{fontenc}
\usepackage[ngerman]{babel}

% ── Fonts ────────────────────────────────────────────────────
\usepackage{lmodern}
\usepackage{microtype}

% ── Page Geometry ────────────────────────────────────────────
\usepackage[
  top=2.8cm,
  bottom=2.8cm,
  left=2.6cm,
  right=2.6cm,
  headheight=28pt
]{geometry}

% ── Colors ───────────────────────────────────────────────────
\usepackage[table]{xcolor}
\definecolor{accent}{RGB}{0, 80, 150}
\definecolor{defblue}{RGB}{220, 235, 255}
\definecolor{thmgreen}{RGB}{220, 245, 225}
\definecolor{examplegray}{RGB}{245, 245, 245}
\definecolor{remarkgold}{RGB}{255, 248, 220}
% table highlight
\definecolor{highlight}{RGB}{255, 195, 215}
% new box colors
\definecolor{satzred-bg}{RGB}{255, 235, 235}
\definecolor{satzred-fr}{RGB}{180, 30, 30}
\definecolor{beweisblue-bg}{RGB}{225, 240, 255}
\definecolor{beweisblue-fr}{RGB}{30, 90, 180}
\definecolor{gedankenpurple-bg}{RGB}{240, 232, 255}
\definecolor{gedankenpurple-fr}{RGB}{110, 50, 180}
\definecolor{exercisegreen-bg}{RGB}{225, 248, 230}
\definecolor{exercisegreen-fr}{RGB}{30, 140, 60}

% ── Header & Footer ──────────────────────────────────────────
\usepackage{fancyhdr}

\fancypagestyle{main}{%
  \fancyhf{}
  \fancyhead[L]{\small\itshape\parbox{0.45\linewidth}{\raggedright\rightmark}}
  \fancyhead[R]{\small\itshape Systemtheorie – Signale und Systeme I}
  \fancyfoot[C]{\small\thepage}
  \renewcommand{\headrulewidth}{0.4pt}
  \renewcommand{\footrulewidth}{0.4pt}
}

% Chapter start pages: use full header too
\renewcommand*{\chapterpagestyle}{main}

% plain style (used for ToC pages)
\fancypagestyle{plain}{%
  \fancyhf{}
  \fancyfoot[C]{\small\thepage}
  \renewcommand{\headrulewidth}{0pt}
  \renewcommand{\footrulewidth}{0.4pt}
}

% ── Math ─────────────────────────────────────────────────────
\usepackage{amsmath}
\usepackage{amssymb}
\usepackage{amsthm}
\usepackage{mathtools}

% ── Graphics & Floats ────────────────────────────────────────
\usepackage{graphicx}
\usepackage{float}
\usepackage{tikz}
\usetikzlibrary{shapes.geometric}
\usepackage{pgfplots}
\pgfplotsset{compat=1.15}
\usepgfplotslibrary{groupplots}

% ── Tables ───────────────────────────────────────────────────
\usepackage{booktabs}
\usepackage{array}
\usepackage{multirow}
\usepackage{makecell}

% ── Lists ────────────────────────────────────────────────────
\usepackage{enumitem}
\setlist[itemize]{itemsep=2pt, topsep=4pt}
\setlist[enumerate]{itemsep=2pt, topsep=4pt}

% ── Colored Boxes ────────────────────────────────────────────
\usepackage[most]{tcolorbox}

\tcbset{
  base style/.style={
    breakable,
    enhanced,
    fonttitle=\bfseries,
    colbacktitle=white,        % weißer Titelbereich
    colframe=black,            % schwarzer Rand
    titlerule=0.4pt,
    attach boxed title to top left={yshift=-2mm, xshift=4mm},
    boxed title style={
      colframe=black,
      colback=white,
      boxrule=0.4pt,
      arc=2pt,
    },
    left=6pt, right=6pt,
    top=6pt, bottom=6pt,
    arc=3pt,
    boxrule=0.6pt,
  }
}

% Definition (blau)
\newtcolorbox{definitionbox}[1][]{
  base style,
  colback=defblue,
  coltitle=accent,
  title={Definition\ifx&#1&\else: #1\fi},
}

% Satz (rot)
\newtcolorbox{satzbox}[1][]{
  base style,
  colback=satzred-bg,
  coltitle=satzred-fr,
  title={Satz\ifx&#1&\else: #1\fi},
}

% Beweis (blau)
\newtcolorbox{beweisbox}[1][]{
  base style,
  colback=beweisblue-bg,
  coltitle=beweisblue-fr,
  title={Beweis\ifx&#1&\else: #1\fi},
}

% Gedankenexperiment (lila)
\newtcolorbox{gedankenbox}[1][]{
  base style,
  colback=gedankenpurple-bg,
  coltitle=gedankenpurple-fr,
  title={Gedankenexperiment\ifx&#1&\else: #1\fi},
}

% Beispiel / Exercise (grün)
\newtcolorbox{exercisebox}[1][]{
  base style,
  colback=exercisegreen-bg,
  coltitle=exercisegreen-fr,
  title={Beispiel\ifx&#1&\else: #1\fi},
}

% Bemerkung (gold)
\newtcolorbox{bemerkungbox}[1][]{
  base style,
  colback=remarkgold,
  coltitle=orange!80!black,
  title={Bemerkung\ifx&#1&\else: #1\fi},
}

% ── Environments (lowercase and capitalized aliases) ──────────
\newenvironment{definition}[1][]{\begin{definitionbox}[#1]}{\end{definitionbox}}
\newenvironment{Definition}[1][]{\begin{definitionbox}[#1]}{\end{definitionbox}}
\newenvironment{satz}[1][]{\begin{satzbox}[#1]}{\end{satzbox}}
\newenvironment{Satz}[1][]{\begin{satzbox}[#1]}{\end{satzbox}}
\newenvironment{beweis}[1][]{\begin{beweisbox}[#1]}{\end{beweisbox}}
\newenvironment{Beweis}[1][]{\begin{beweisbox}[#1]}{\end{beweisbox}}
\newenvironment{gedankenexperiment}[1][]{\begin{gedankenbox}[#1]}{\end{gedankenbox}}
\newenvironment{Gedankenexperiment}[1][]{\begin{gedankenbox}[#1]}{\end{gedankenbox}}
\newenvironment{beispiel}[1][]{\begin{exercisebox}[#1]}{\end{exercisebox}}
\newenvironment{Beispiel}[1][]{\begin{exercisebox}[#1]}{\end{exercisebox}}
\newenvironment{bemerkung}[1][]{\begin{bemerkungbox}[#1]}{\end{bemerkungbox}}
\newenvironment{Bemerkung}[1][]{\begin{bemerkungbox}[#1]}{\end{bemerkungbox}}

% ── Heading Spacing ──────────────────────────────────────────
% Reduce the space below chapter, section, subsection headings
\RedeclareSectionCommand[beforeskip=0.8\baselineskip, afterskip=0.5\baselineskip]{chapter}
\RedeclareSectionCommand[afterskip=0.3\baselineskip]{section}
\RedeclareSectionCommand[afterskip=0.2\baselineskip]{subsection}

% ── Hyperlinks ───────────────────────────────────────────────
\usepackage[
  colorlinks=true,
  linkcolor=accent,
  citecolor=accent,
  urlcolor=accent,
  bookmarksnumbered=true,
  pdfauthor={Philipp Lehmann},
  pdftitle={Systemtheorie – Signale und Systeme I},
]{hyperref}

% ============================================================

% ── Encoding & Language ──────────────────────────────────────
\usepackage[utf8]{inputenc}
\usepackage[T1]{fontenc}
\usepackage[ngerman]{babel}

% ── Fonts ────────────────────────────────────────────────────
\usepackage{lmodern}
\usepackage{microtype}

% ── Page Geometry ────────────────────────────────────────────
\usepackage[
  top=2.8cm,
  bottom=2.8cm,
  left=2.6cm,
  right=2.6cm,
  headheight=28pt
]{geometry}

% ── Colors ───────────────────────────────────────────────────
\usepackage[table]{xcolor}
\definecolor{accent}{RGB}{0, 80, 150}
\definecolor{defblue}{RGB}{220, 235, 255}
\definecolor{thmgreen}{RGB}{220, 245, 225}
\definecolor{examplegray}{RGB}{245, 245, 245}
\definecolor{remarkgold}{RGB}{255, 248, 220}
% table highlight
\definecolor{highlight}{RGB}{255, 195, 215}
% new box colors
\definecolor{satzred-bg}{RGB}{255, 235, 235}
\definecolor{satzred-fr}{RGB}{180, 30, 30}
\definecolor{beweisblue-bg}{RGB}{225, 240, 255}
\definecolor{beweisblue-fr}{RGB}{30, 90, 180}
\definecolor{gedankenpurple-bg}{RGB}{240, 232, 255}
\definecolor{gedankenpurple-fr}{RGB}{110, 50, 180}
\definecolor{exercisegreen-bg}{RGB}{225, 248, 230}
\definecolor{exercisegreen-fr}{RGB}{30, 140, 60}

% ── Header & Footer ──────────────────────────────────────────
\usepackage{fancyhdr}

\fancypagestyle{main}{%
  \fancyhf{}
  \fancyhead[L]{\small\itshape\parbox{0.45\linewidth}{\raggedright\rightmark}}
  \fancyhead[R]{\small\itshape Systemtheorie – Signale und Systeme I}
  \fancyfoot[C]{\small\thepage}
  \renewcommand{\headrulewidth}{0.4pt}
  \renewcommand{\footrulewidth}{0.4pt}
}

% Chapter start pages: use full header too
\renewcommand*{\chapterpagestyle}{main}

% plain style (used for ToC pages)
\fancypagestyle{plain}{%
  \fancyhf{}
  \fancyfoot[C]{\small\thepage}
  \renewcommand{\headrulewidth}{0pt}
  \renewcommand{\footrulewidth}{0.4pt}
}

% ── Math ─────────────────────────────────────────────────────
\usepackage{amsmath}
\usepackage{amssymb}
\usepackage{amsthm}
\usepackage{mathtools}

% ── Graphics & Floats ────────────────────────────────────────
\usepackage{graphicx}
\usepackage{float}
\usepackage{tikz}
\usetikzlibrary{shapes.geometric}
\usepackage{pgfplots}
\pgfplotsset{compat=1.15}
\usepgfplotslibrary{groupplots}

% ── Tables ───────────────────────────────────────────────────
\usepackage{booktabs}
\usepackage{array}
\usepackage{multirow}
\usepackage{makecell}

% ── Lists ────────────────────────────────────────────────────
\usepackage{enumitem}
\setlist[itemize]{itemsep=2pt, topsep=4pt}
\setlist[enumerate]{itemsep=2pt, topsep=4pt}

% ── Colored Boxes ────────────────────────────────────────────
\usepackage[most]{tcolorbox}

\tcbset{
  base style/.style={
    breakable,
    enhanced,
    fonttitle=\bfseries,
    colbacktitle=white,        % weißer Titelbereich
    colframe=black,            % schwarzer Rand
    titlerule=0.4pt,
    attach boxed title to top left={yshift=-2mm, xshift=4mm},
    boxed title style={
      colframe=black,
      colback=white,
      boxrule=0.4pt,
      arc=2pt,
    },
    left=6pt, right=6pt,
    top=6pt, bottom=6pt,
    arc=3pt,
    boxrule=0.6pt,
  }
}

% Definition (blau)
\newtcolorbox{definitionbox}[1][]{
  base style,
  colback=defblue,
  coltitle=accent,
  title={Definition\ifx&#1&\else: #1\fi},
}

% Satz (rot)
\newtcolorbox{satzbox}[1][]{
  base style,
  colback=satzred-bg,
  coltitle=satzred-fr,
  title={Satz\ifx&#1&\else: #1\fi},
}

% Beweis (blau)
\newtcolorbox{beweisbox}[1][]{
  base style,
  colback=beweisblue-bg,
  coltitle=beweisblue-fr,
  title={Beweis\ifx&#1&\else: #1\fi},
}

% Gedankenexperiment (lila)
\newtcolorbox{gedankenbox}[1][]{
  base style,
  colback=gedankenpurple-bg,
  coltitle=gedankenpurple-fr,
  title={Gedankenexperiment\ifx&#1&\else: #1\fi},
}

% Beispiel / Exercise (grün)
\newtcolorbox{exercisebox}[1][]{
  base style,
  colback=exercisegreen-bg,
  coltitle=exercisegreen-fr,
  title={Beispiel\ifx&#1&\else: #1\fi},
}

% Bemerkung (gold)
\newtcolorbox{bemerkungbox}[1][]{
  base style,
  colback=remarkgold,
  coltitle=orange!80!black,
  title={Bemerkung\ifx&#1&\else: #1\fi},
}

% ── Environments (lowercase and capitalized aliases) ──────────
\newenvironment{definition}[1][]{\begin{definitionbox}[#1]}{\end{definitionbox}}
\newenvironment{Definition}[1][]{\begin{definitionbox}[#1]}{\end{definitionbox}}
\newenvironment{satz}[1][]{\begin{satzbox}[#1]}{\end{satzbox}}
\newenvironment{Satz}[1][]{\begin{satzbox}[#1]}{\end{satzbox}}
\newenvironment{beweis}[1][]{\begin{beweisbox}[#1]}{\end{beweisbox}}
\newenvironment{Beweis}[1][]{\begin{beweisbox}[#1]}{\end{beweisbox}}
\newenvironment{gedankenexperiment}[1][]{\begin{gedankenbox}[#1]}{\end{gedankenbox}}
\newenvironment{Gedankenexperiment}[1][]{\begin{gedankenbox}[#1]}{\end{gedankenbox}}
\newenvironment{beispiel}[1][]{\begin{exercisebox}[#1]}{\end{exercisebox}}
\newenvironment{Beispiel}[1][]{\begin{exercisebox}[#1]}{\end{exercisebox}}
\newenvironment{bemerkung}[1][]{\begin{bemerkungbox}[#1]}{\end{bemerkungbox}}
\newenvironment{Bemerkung}[1][]{\begin{bemerkungbox}[#1]}{\end{bemerkungbox}}

% ── Heading Spacing ──────────────────────────────────────────
% Reduce the space below chapter, section, subsection headings
\RedeclareSectionCommand[beforeskip=0.8\baselineskip, afterskip=0.5\baselineskip]{chapter}
\RedeclareSectionCommand[afterskip=0.3\baselineskip]{section}
\RedeclareSectionCommand[afterskip=0.2\baselineskip]{subsection}

% ── Hyperlinks ───────────────────────────────────────────────
\usepackage[
  colorlinks=true,
  linkcolor=accent,
  citecolor=accent,
  urlcolor=accent,
  bookmarksnumbered=true,
  pdfauthor={Philipp Lehmann},
  pdftitle={Systemtheorie – Signale und Systeme I},
]{hyperref}

% ============================================================

% ── Encoding & Language ──────────────────────────────────────
\usepackage[utf8]{inputenc}
\usepackage[T1]{fontenc}
\usepackage[ngerman]{babel}

% ── Fonts ────────────────────────────────────────────────────
\usepackage{lmodern}
\usepackage{microtype}

% ── Page Geometry ────────────────────────────────────────────
\usepackage[
  top=2.8cm,
  bottom=2.8cm,
  left=2.6cm,
  right=2.6cm,
  headheight=28pt
]{geometry}

% ── Colors ───────────────────────────────────────────────────
\usepackage[table]{xcolor}
\definecolor{accent}{RGB}{0, 80, 150}
\definecolor{defblue}{RGB}{220, 235, 255}
\definecolor{thmgreen}{RGB}{220, 245, 225}
\definecolor{examplegray}{RGB}{245, 245, 245}
\definecolor{remarkgold}{RGB}{255, 248, 220}
% table highlight
\definecolor{highlight}{RGB}{255, 195, 215}
% new box colors
\definecolor{satzred-bg}{RGB}{255, 235, 235}
\definecolor{satzred-fr}{RGB}{180, 30, 30}
\definecolor{beweisblue-bg}{RGB}{225, 240, 255}
\definecolor{beweisblue-fr}{RGB}{30, 90, 180}
\definecolor{gedankenpurple-bg}{RGB}{240, 232, 255}
\definecolor{gedankenpurple-fr}{RGB}{110, 50, 180}
\definecolor{exercisegreen-bg}{RGB}{225, 248, 230}
\definecolor{exercisegreen-fr}{RGB}{30, 140, 60}

% ── Header & Footer ──────────────────────────────────────────
\usepackage{fancyhdr}

\fancypagestyle{main}{%
  \fancyhf{}
  \fancyhead[L]{\small\itshape\parbox{0.45\linewidth}{\raggedright\rightmark}}
  \fancyhead[R]{\small\itshape Systemtheorie – Signale und Systeme I}
  \fancyfoot[C]{\small\thepage}
  \renewcommand{\headrulewidth}{0.4pt}
  \renewcommand{\footrulewidth}{0.4pt}
}

% Chapter start pages: use full header too
\renewcommand*{\chapterpagestyle}{main}

% plain style (used for ToC pages)
\fancypagestyle{plain}{%
  \fancyhf{}
  \fancyfoot[C]{\small\thepage}
  \renewcommand{\headrulewidth}{0pt}
  \renewcommand{\footrulewidth}{0.4pt}
}

% ── Math ─────────────────────────────────────────────────────
\usepackage{amsmath}
\usepackage{amssymb}
\usepackage{amsthm}
\usepackage{mathtools}

% ── Graphics & Floats ────────────────────────────────────────
\usepackage{graphicx}
\usepackage{float}
\usepackage{tikz}
\usetikzlibrary{shapes.geometric}
\usepackage{pgfplots}
\pgfplotsset{compat=1.15}
\usepgfplotslibrary{groupplots}

% ── Tables ───────────────────────────────────────────────────
\usepackage{booktabs}
\usepackage{array}
\usepackage{multirow}
\usepackage{makecell}

% ── Lists ────────────────────────────────────────────────────
\usepackage{enumitem}
\setlist[itemize]{itemsep=2pt, topsep=4pt}
\setlist[enumerate]{itemsep=2pt, topsep=4pt}

% ── Colored Boxes ────────────────────────────────────────────
\usepackage[most]{tcolorbox}

\tcbset{
  base style/.style={
    breakable,
    enhanced,
    fonttitle=\bfseries,
    colbacktitle=white,        % weißer Titelbereich
    colframe=black,            % schwarzer Rand
    titlerule=0.4pt,
    attach boxed title to top left={yshift=-2mm, xshift=4mm},
    boxed title style={
      colframe=black,
      colback=white,
      boxrule=0.4pt,
      arc=2pt,
    },
    left=6pt, right=6pt,
    top=6pt, bottom=6pt,
    arc=3pt,
    boxrule=0.6pt,
  }
}

% Definition (blau)
\newtcolorbox{definitionbox}[1][]{
  base style,
  colback=defblue,
  coltitle=accent,
  title={Definition\ifx&#1&\else: #1\fi},
}

% Satz (rot)
\newtcolorbox{satzbox}[1][]{
  base style,
  colback=satzred-bg,
  coltitle=satzred-fr,
  title={Satz\ifx&#1&\else: #1\fi},
}

% Beweis (blau)
\newtcolorbox{beweisbox}[1][]{
  base style,
  colback=beweisblue-bg,
  coltitle=beweisblue-fr,
  title={Beweis\ifx&#1&\else: #1\fi},
}

% Gedankenexperiment (lila)
\newtcolorbox{gedankenbox}[1][]{
  base style,
  colback=gedankenpurple-bg,
  coltitle=gedankenpurple-fr,
  title={Gedankenexperiment\ifx&#1&\else: #1\fi},
}

% Beispiel / Exercise (grün)
\newtcolorbox{exercisebox}[1][]{
  base style,
  colback=exercisegreen-bg,
  coltitle=exercisegreen-fr,
  title={Beispiel\ifx&#1&\else: #1\fi},
}

% Bemerkung (gold)
\newtcolorbox{bemerkungbox}[1][]{
  base style,
  colback=remarkgold,
  coltitle=orange!80!black,
  title={Bemerkung\ifx&#1&\else: #1\fi},
}

% ── Environments (lowercase and capitalized aliases) ──────────
\newenvironment{definition}[1][]{\begin{definitionbox}[#1]}{\end{definitionbox}}
\newenvironment{Definition}[1][]{\begin{definitionbox}[#1]}{\end{definitionbox}}
\newenvironment{satz}[1][]{\begin{satzbox}[#1]}{\end{satzbox}}
\newenvironment{Satz}[1][]{\begin{satzbox}[#1]}{\end{satzbox}}
\newenvironment{beweis}[1][]{\begin{beweisbox}[#1]}{\end{beweisbox}}
\newenvironment{Beweis}[1][]{\begin{beweisbox}[#1]}{\end{beweisbox}}
\newenvironment{gedankenexperiment}[1][]{\begin{gedankenbox}[#1]}{\end{gedankenbox}}
\newenvironment{Gedankenexperiment}[1][]{\begin{gedankenbox}[#1]}{\end{gedankenbox}}
\newenvironment{beispiel}[1][]{\begin{exercisebox}[#1]}{\end{exercisebox}}
\newenvironment{Beispiel}[1][]{\begin{exercisebox}[#1]}{\end{exercisebox}}
\newenvironment{bemerkung}[1][]{\begin{bemerkungbox}[#1]}{\end{bemerkungbox}}
\newenvironment{Bemerkung}[1][]{\begin{bemerkungbox}[#1]}{\end{bemerkungbox}}

% ── Heading Spacing ──────────────────────────────────────────
% Reduce the space below chapter, section, subsection headings
\RedeclareSectionCommand[beforeskip=0.8\baselineskip, afterskip=0.5\baselineskip]{chapter}
\RedeclareSectionCommand[afterskip=0.3\baselineskip]{section}
\RedeclareSectionCommand[afterskip=0.2\baselineskip]{subsection}

% ── Hyperlinks ───────────────────────────────────────────────
\usepackage[
  colorlinks=true,
  linkcolor=accent,
  citecolor=accent,
  urlcolor=accent,
  bookmarksnumbered=true,
  pdfauthor={Philipp Lehmann},
  pdftitle={Systemtheorie – Signale und Systeme I},
]{hyperref}

% ============================================================

% ── Encoding & Language ──────────────────────────────────────
\usepackage[utf8]{inputenc}
\usepackage[T1]{fontenc}
\usepackage[ngerman]{babel}

% ── Fonts ────────────────────────────────────────────────────
\usepackage{lmodern}
\usepackage{microtype}

% ── Page Geometry ────────────────────────────────────────────
\usepackage[
  top=2.8cm,
  bottom=2.8cm,
  left=2.6cm,
  right=2.6cm,
  headheight=28pt
]{geometry}

% ── Colors ───────────────────────────────────────────────────
\usepackage[table]{xcolor}
\definecolor{accent}{RGB}{0, 80, 150}
\definecolor{defblue}{RGB}{220, 235, 255}
\definecolor{thmgreen}{RGB}{220, 245, 225}
\definecolor{examplegray}{RGB}{245, 245, 245}
\definecolor{remarkgold}{RGB}{255, 248, 220}
% table highlight
\definecolor{highlight}{RGB}{255, 195, 215}
% new box colors
\definecolor{satzred-bg}{RGB}{255, 235, 235}
\definecolor{satzred-fr}{RGB}{180, 30, 30}
\definecolor{beweisblue-bg}{RGB}{225, 240, 255}
\definecolor{beweisblue-fr}{RGB}{30, 90, 180}
\definecolor{gedankenpurple-bg}{RGB}{240, 232, 255}
\definecolor{gedankenpurple-fr}{RGB}{110, 50, 180}
\definecolor{exercisegreen-bg}{RGB}{225, 248, 230}
\definecolor{exercisegreen-fr}{RGB}{30, 140, 60}

% ── Header & Footer ──────────────────────────────────────────
\usepackage{fancyhdr}

\fancypagestyle{main}{%
  \fancyhf{}
  \fancyhead[L]{\small\itshape\parbox{0.45\linewidth}{\raggedright\rightmark}}
  \fancyhead[R]{\small\itshape Systemtheorie – Signale und Systeme I}
  \fancyfoot[C]{\small\thepage}
  \renewcommand{\headrulewidth}{0.4pt}
  \renewcommand{\footrulewidth}{0.4pt}
}

% Chapter start pages: use full header too
\renewcommand*{\chapterpagestyle}{main}

% plain style (used for ToC pages)
\fancypagestyle{plain}{%
  \fancyhf{}
  \fancyfoot[C]{\small\thepage}
  \renewcommand{\headrulewidth}{0pt}
  \renewcommand{\footrulewidth}{0.4pt}
}

% ── Math ─────────────────────────────────────────────────────
\usepackage{amsmath}
\usepackage{amssymb}
\usepackage{amsthm}
\usepackage{mathtools}

% ── Graphics & Floats ────────────────────────────────────────
\usepackage{graphicx}
\usepackage{float}
\usepackage{tikz}
\usetikzlibrary{shapes.geometric}
\usepackage{pgfplots}
\pgfplotsset{compat=1.15}
\usepgfplotslibrary{groupplots}

% ── Tables ───────────────────────────────────────────────────
\usepackage{booktabs}
\usepackage{array}
\usepackage{multirow}
\usepackage{makecell}

% ── Lists ────────────────────────────────────────────────────
\usepackage{enumitem}
\setlist[itemize]{itemsep=2pt, topsep=4pt}
\setlist[enumerate]{itemsep=2pt, topsep=4pt}

% ── Colored Boxes ────────────────────────────────────────────
\usepackage[most]{tcolorbox}

\tcbset{
  base style/.style={
    breakable,
    enhanced,
    fonttitle=\bfseries,
    colbacktitle=white,        % weißer Titelbereich
    colframe=black,            % schwarzer Rand
    titlerule=0.4pt,
    attach boxed title to top left={yshift=-2mm, xshift=4mm},
    boxed title style={
      colframe=black,
      colback=white,
      boxrule=0.4pt,
      arc=2pt,
    },
    left=6pt, right=6pt,
    top=6pt, bottom=6pt,
    arc=3pt,
    boxrule=0.6pt,
  }
}

% Definition (blau)
\newtcolorbox{definitionbox}[1][]{
  base style,
  colback=defblue,
  coltitle=accent,
  title={Definition\ifx&#1&\else: #1\fi},
}

% Satz (rot)
\newtcolorbox{satzbox}[1][]{
  base style,
  colback=satzred-bg,
  coltitle=satzred-fr,
  title={Satz\ifx&#1&\else: #1\fi},
}

% Beweis (blau)
\newtcolorbox{beweisbox}[1][]{
  base style,
  colback=beweisblue-bg,
  coltitle=beweisblue-fr,
  title={Beweis\ifx&#1&\else: #1\fi},
}

% Gedankenexperiment (lila)
\newtcolorbox{gedankenbox}[1][]{
  base style,
  colback=gedankenpurple-bg,
  coltitle=gedankenpurple-fr,
  title={Gedankenexperiment\ifx&#1&\else: #1\fi},
}

% Beispiel / Exercise (grün)
\newtcolorbox{exercisebox}[1][]{
  base style,
  colback=exercisegreen-bg,
  coltitle=exercisegreen-fr,
  title={Beispiel\ifx&#1&\else: #1\fi},
}

% Bemerkung (gold)
\newtcolorbox{bemerkungbox}[1][]{
  base style,
  colback=remarkgold,
  coltitle=orange!80!black,
  title={Bemerkung\ifx&#1&\else: #1\fi},
}

% ── Environments (lowercase and capitalized aliases) ──────────
\newenvironment{definition}[1][]{\begin{definitionbox}[#1]}{\end{definitionbox}}
\newenvironment{Definition}[1][]{\begin{definitionbox}[#1]}{\end{definitionbox}}
\newenvironment{satz}[1][]{\begin{satzbox}[#1]}{\end{satzbox}}
\newenvironment{Satz}[1][]{\begin{satzbox}[#1]}{\end{satzbox}}
\newenvironment{beweis}[1][]{\begin{beweisbox}[#1]}{\end{beweisbox}}
\newenvironment{Beweis}[1][]{\begin{beweisbox}[#1]}{\end{beweisbox}}
\newenvironment{gedankenexperiment}[1][]{\begin{gedankenbox}[#1]}{\end{gedankenbox}}
\newenvironment{Gedankenexperiment}[1][]{\begin{gedankenbox}[#1]}{\end{gedankenbox}}
\newenvironment{beispiel}[1][]{\begin{exercisebox}[#1]}{\end{exercisebox}}
\newenvironment{Beispiel}[1][]{\begin{exercisebox}[#1]}{\end{exercisebox}}
\newenvironment{bemerkung}[1][]{\begin{bemerkungbox}[#1]}{\end{bemerkungbox}}
\newenvironment{Bemerkung}[1][]{\begin{bemerkungbox}[#1]}{\end{bemerkungbox}}

% ── Heading Spacing ──────────────────────────────────────────
% Reduce the space below chapter, section, subsection headings
\RedeclareSectionCommand[beforeskip=0.8\baselineskip, afterskip=0.5\baselineskip]{chapter}
\RedeclareSectionCommand[afterskip=0.3\baselineskip]{section}
\RedeclareSectionCommand[afterskip=0.2\baselineskip]{subsection}

% ── Hyperlinks ───────────────────────────────────────────────
\usepackage[
  colorlinks=true,
  linkcolor=accent,
  citecolor=accent,
  urlcolor=accent,
  bookmarksnumbered=true,
  pdfauthor={Philipp Lehmann},
  pdftitle={Systemtheorie – Signale und Systeme I},
]{hyperref}
